\section{Shortest Common Supersequence}
\label{sec:dp-shortest-common-supersequence}

\problemstatement{Given two strings $s_1$ (length $n$) and $s_2$ (length
$m$), construct the \textbf{shortest} string that has \textbf{both}
$s_1$ and $s_2$ as subsequences.}

\begin{patternbox}
\emph{``Shortest string containing both as subsequences''} is the
``opposite'' framing of Section~\ref{sec:dp-lcs}'s LCS, and the two are
linked by a clean formula: every character the two strings \emph{share}
in their LCS needs to appear only \textbf{once} in the supersequence
(since it serves both strings' subsequence requirement simultaneously),
while every other character of each string must still appear --
collapsing the problem into one already-solved LCS computation plus a
reconstruction walk almost identical to Section~\ref{sec:dp-print-lcs}'s.
\end{patternbox}

\subsection*{Deriving the length formula}

If we simply concatenated $s_1$ and $s_2$, the result trivially contains
both as subsequences, with length $n+m$ -- but every character that is
part of \emph{some} common subsequence is being counted \textbf{twice}
in this concatenation (once from each string), when it only needs to
appear once. The longest common subsequence is exactly the largest set
of characters we can ``merge'' this way, so:
\[
  |\textit{SCS}| = n + m - |\textit{LCS}(s_1, s_2)|.
\]

\subsection*{Constructing the actual supersequence}

Run the same backward walk as Section~\ref{sec:dp-print-lcs}, but
instead of recording \emph{only} matched characters, record
\textbf{every} character visited along the way: on a match, record the
(shared) character once and move diagonally; on a mismatch, record
whichever of $s_1[i-1]$ or $s_2[j-1]$ corresponds to the direction moved
(following the same $\textit{dp}[i-1][j]$ vs.\ $\textit{dp}[i][j-1]$
comparison as before), and move only that one pointer. Once either
pointer reaches $0$, append the entirety of whatever remains of the
other string (every leftover character must still appear, in order,
since that string's subsequence requirement is otherwise unmet).

\subsection*{Algorithm}

\begin{enumerate}
  \item Build the LCS tabulation table for $s_1, s_2$ (as in
        Section~\ref{sec:dp-lcs}).
  \item $i \gets n$, $j \gets m$, $\textit{result} \gets \texttt{""}$.
  \item While $i>0$ and $j>0$:
  \begin{enumerate}
    \item If $s_1[i-1]=s_2[j-1]$: prepend $s_1[i-1]$; $i{-}{-}$,
          $j{-}{-}$.
    \item Else if $\textit{dp}[i-1][j] \ge \textit{dp}[i][j-1]$:
          prepend $s_1[i-1]$; $i{-}{-}$.
    \item Else: prepend $s_2[j-1]$; $j{-}{-}$.
  \end{enumerate}
  \item Prepend any remaining prefix of $s_1$ (if $i>0$) or $s_2$ (if
        $j>0$).
  \item Return \textit{result}.
\end{enumerate}

\subsection*{Dry run}

\dryrun{Take $s_1=\texttt{"abac"}$, $s_2=\texttt{"cab"}$.}

LCS of \texttt{"abac"} and \texttt{"cab"} is \texttt{"ab"} (length $2$),
so $|\textit{SCS}| = 4+3-2 = 5$. Constructing: walking backward through
the LCS table merges the two strings, interleaving non-matched
characters with the two shared \texttt{a} and \texttt{b}, producing one
valid result such as \texttt{"cabac"} (length $5$) -- which indeed
contains \texttt{"abac"} as a subsequence (positions $1,3,4,5$... more
directly, removing the leading \texttt{c} from \texttt{"cabac"} leaves
\texttt{"abac"} exactly) and \texttt{"cab"} as a subsequence (the first
three characters).

\subsection*{C++ implementation}

\begin{lstlisting}
#include <bits/stdc++.h>
using namespace std;

string shortestCommonSupersequence(string& s1, string& s2) {
    int n = s1.size(), m = s2.size();
    vector<vector<int>> dp(n + 1, vector<int>(m + 1, 0));

    for (int i = 1; i <= n; i++) {
        for (int j = 1; j <= m; j++) {
            if (s1[i - 1] == s2[j - 1]) {
                dp[i][j] = 1 + dp[i - 1][j - 1];
            } else {
                dp[i][j] = max(dp[i - 1][j], dp[i][j - 1]);
            }
        }
    }

    string result;
    int i = n, j = m;
    while (i > 0 && j > 0) {
        if (s1[i - 1] == s2[j - 1]) {
            result.push_back(s1[i - 1]);
            i--; j--;
        } else if (dp[i - 1][j] >= dp[i][j - 1]) {
            result.push_back(s1[i - 1]);
            i--;
        } else {
            result.push_back(s2[j - 1]);
            j--;
        }
    }
    while (i > 0) { result.push_back(s1[i - 1]); i--; }
    while (j > 0) { result.push_back(s2[j - 1]); j--; }

    reverse(result.begin(), result.end());
    return result;
}

int main() {
    string s1 = "abac", s2 = "cab";
    string scs = shortestCommonSupersequence(s1, s2);
    cout << scs << " (length " << scs.size() << ")" << endl;
    return 0;
}
\end{lstlisting}

\begin{complexitybox}
\textbf{Time Complexity.} $O(n \cdot m)$ to build the LCS table, plus
$O(n+m)$ for the reconstruction walk.
\textbf{Space Complexity.} $O(n \cdot m)$ for the full table (needed for
reconstruction, as in Section~\ref{sec:dp-print-lcs}).
\end{complexitybox}

\begin{takeawaybox}
Shortest Common Supersequence is not a new DP recurrence at all -- it is
LCS's length formula, $n+m-|\textit{LCS}|$, plus a reconstruction walk
that is a small variation on Print LCS's (recording every visited
character, not only the matches). Whenever a problem can be phrased as
``the complement or merge of an already-solved quantity,'' check the
arithmetic relationship before reaching for a new table.
\end{takeawaybox}
